\documentclass[12pt,a4paper]{article}

% --- PAQUETES DE CONFIGURACIÓN ---
\usepackage[spanish]{babel}
\usepackage[utf8]{inputenc}
\usepackage[T1]{fontenc}
\usepackage{geometry}
\geometry{margin=2.5cm}
\setlength{\emergencystretch}{1em} % margen extra para evitar desbordes
\usepackage{listings}
\usepackage{xcolor}
\usepackage{endnotes}
\usepackage{amsmath}
\usepackage{seqsplit}
\usepackage[numbers]{natbib}
\usepackage{hyperref} % hyperref se carga en último lugar

% --- CONFIGURACIÓN DE APARIENCIA PARA CÓDIGO ---
\definecolor{codegreen}{rgb}{0,0.6,0}
\definecolor{codegray}{rgb}{0.5,0.5,0.5}
\definecolor{codeblue}{rgb}{0,0,0.8}
\definecolor{backcolour}{rgb}{0.96,0.96,0.96}

\lstdefinestyle{mystyle}{
    backgroundcolor=\color{backcolour},
    commentstyle=\color{codegreen},
    keywordstyle=\color{codeblue},
    numberstyle=\tiny\color{codegray},
    stringstyle=\color{red},
    basicstyle=\ttfamily\small,
    breakatwhitespace=false,
    breaklines=true,
    captionpos=b,
    keepspaces=true,
    numbers=left,
    numbersep=5pt,
    showspaces=false,
    showstringspaces=false,
    showtabs=false,
    tabsize=2,
    literate=
        {á}{{\'a}}1
        {é}{{\'e}}1
        {í}{{\'i}}1
        {ó}{{\'o}}1
        {ú}{{\'u}}1
        {ñ}{{\~n}}1
}
\lstset{style=mystyle}

\lstdefinelanguage{properties}{
  sensitive=false,
  morecomment=[l]{\#},
  morecomment=[l]{!}
}

% --- DATOS DEL DOCUMENTO ---
\title{Bases de Datos Relacionales de Alto Rendimiento: Amazon Aurora}
\date{\normalsize\today}

\input{datos_proyecto}

\begin{document}

\maketitle

\begin{abstract}
Este documento presenta la arquitectura de Amazon Aurora, un motor relacional
compatible con MySQL y PostgreSQL que separa la capa de cómputo de una capa de
almacenamiento distribuida y replicada. Tras describir su modelo de
almacenamiento, sus mecanismos de alta disponibilidad y sus opciones de
seguridad, se detalla su aplicación en \proyecto: el despliegue automatizado de
un clúster Aurora PostgreSQL mediante la CLI de AWS, la gestión de la
contraseña maestra con AWS Secrets Manager sin que esta transite por la línea
de comandos y la integración del clúster con un \textit{backend} Spring Boot.
\end{abstract}

{
\let\clearpage\relax
\tableofcontents
}

\section{Introducción}

Conectar una aplicación Spring Boot a una base de datos relacional no es nada nuevo: hasta ahora \proyecto~usaba H2, y pasar a MySQL o PostgreSQL consistiría en añadir el driver JDBC y cambiar la URL, el usuario y la contraseña. Aurora no es distinta en ese aspecto. Si nos limitamos a crear un clúster con una instancia y a apuntar la URL JDBC a su \textit{endpoint}, Aurora es una base de datos más: el código Java de la aplicación (entidades, repositorios y servicios) no cambia respecto a H2.

La diferencia aparece cuando la aplicación tiene que cumplir requisitos de tolerancia a fallos, disponibilidad y escalabilidad. Con un poco de configuración, sin programar nada nuevo, Aurora mantiene réplicas que sustituyen a la instancia de escritura si esta falla, sin pérdida de datos, y distribuye automáticamente las lecturas (no las escrituras) entre las réplicas, lo que permite atender más peticiones simultáneamente.

\section{Amazon Aurora}\label{sec:aurora}

Amazon Aurora es un motor de base de datos relacional compatible con MySQL y PostgreSQL diseñado específicamente para la nube. Según las pruebas realizadas por el propio fabricante con \textit{benchmarks} estándar como SysBench, su rendimiento llega a multiplicar por seis el de MySQL y PostgreSQL estándar sobre un hardware similar \cite{aws_aurora_performance}, si bien estas cifras dependen en gran medida de la carga de trabajo y de la configuración empleada.

En una base de datos relacional convencional, cada servidor ejecuta el motor de base de datos y guarda los datos en su propio disco. Aurora separa esas dos funciones en dos capas:

\begin{itemize}
    \item \textbf{La capa de cómputo} la forman las \textbf{instancias} (\textit{DB instances}; AWS también las llama \textit{database nodes}). Cada instancia es un servidor que ejecuta el motor de base de datos y atiende las consultas, pero no guarda los datos en su disco. Hay dos papeles: la \textbf{instancia de escritura} (\textit{writer}), que es única y la única que acepta escrituras (así como lecturas), y las \textbf{réplicas de lectura} (\textit{readers}), hasta quince, que solo atienden lecturas. 
    
    \item \textbf{La capa de almacenamiento} es un servicio distribuido que AWS gestiona y que guarda los datos de las instancias. Los datos no están en el disco de ninguna instancia: se almacenan en un volumen virtual y autoescalable, uno por clúster (el \textit{cluster volume} de la sección~\ref{sec:despliegue}), que comparten todas sus instancias. Los datos se dividen en segmentos de 10 GB, y el volumen crece automáticamente añadiendo segmentos según la demanda. Cada segmento se guarda seis veces, dos en cada una de tres zonas de disponibilidad de la región. Cuáles son esas tres zonas lo decide AWS, y es transparente para el usuario. Esto es, precisamente, la característica que proporciona \textbf{tolerancia a fallos} a este elemento arquitectónico. Esta replicación \textbf{no depende de las subredes ni de las zonas de disponibilidad} en las que se coloquen las instancias.
    
\end{itemize}

Por eso, en Aurora una réplica de lectura no es una copia de los datos. En la replicación clásica de PostgreSQL o MySQL, cada réplica tiene su propio disco con una copia completa de la base de datos, que mantiene al día aplicando los cambios del servidor principal. En Aurora, la réplica no guarda datos propios: es una instancia más conectada al mismo volumen. En consecuencia, añadir una réplica no duplica el almacenamiento ni obliga a copiar datos, y la durabilidad de los datos no depende de ella, sino de las seis copias de la capa de almacenamiento. Lo que aporta una réplica es cómputo: atiende lecturas y está preparada para sustituir a la instancia de escritura si esta falla.

\section{Despliegue de Amazon Aurora PostgreSQL}\label{sec:despliegue}

Un \textbf{clúster} (\textit{DB cluster}) es la unidad que agrupa todo lo necesario para una base de datos:
\begin{itemize}
    \item \textbf{Un volumen de almacenamiento} (\textit{cluster volume}), único y compartido por todas las instancias del clúster, donde están los datos. Se guarda en la capa de almacenamiento que hemos mencionado en la sección anterior.
    
    \item \textbf{Las instancias}: una de escritura y, opcionalmente, réplicas de lectura, todas conectadas al mismo volumen de almacenamiento. 
    
    \item \textbf{La configuración común}: motor y versión, usuario \textit{master}, base de datos inicial, puerto, grupo de seguridad, \textit{subnet group}, copias de seguridad y los \textit{endpoints} (las direcciones que aparecen en \texttt{jdbc:aws-wrapper:postgresql://...}) por los que se accede a la base de datos.
\end{itemize}

Por este motivo, el clúster y las instancias se crean por separado: primero \texttt{\seqsplit{create-db-cluster}} crea el volumen y la configuración común, y cada \texttt{\seqsplit{create-db-instance}} añade una instancia. 

\subsection{Creación del clúster}\label{sec:cluster}

Un clúster debe estar asociado a, como mínimo, dos subredes en dos zonas de disponibilidad distintas. Esto es necesario para ofrecer una adecuada disponibilidad del servicio, como se indicará en la sección~\ref{sec:writer}.

Aurora no busca las subredes por su cuenta, ni crea un grupo de seguridad propio: usa los que ya creó \texttt{ec2.sh}. Los scripts de despliegue comparten un fichero de estado, \texttt{aws-scripts/lab-state.json}, en el que cada \texttt{create} va guardando los identificadores de lo que crea. Antes de Aurora, \texttt{ec2.sh} ya ha escrito, entre otras, estas claves:

\begin{itemize}
    \item \texttt{GroupId}: el grupo de seguridad \texttt{eventhub-sg} de la instancia EC2.
    \item \texttt{SubnetIds}: las dos primeras subredes de la VPC por defecto de AWS Academy. La VPC por defecto tiene una subred por zona de disponibilidad, así que las dos están en zonas distintas.
\end{itemize}

\noindent \texttt{aurora.sh} las lee con la función \texttt{state\_require} de \texttt{jq-functions.sh}:

\begin{lstlisting}[language=bash,basicstyle=\ttfamily\footnotesize]
SUBNET_IDS=$(state_require SubnetIds)
SG_ID=$(state_require GroupId)
\end{lstlisting}

\vspace{1em}

\noindent \texttt{state\_require} devuelve el valor de la clave, o termina con error si no existe, lo que ocurre si no se ha ejecutado antes \texttt{ec2.sh create}. \texttt{SubnetIds} es un array JSON, y \texttt{state\_require} lo devuelve en una sola línea, con los identificadores separados por espacios.

El comando \texttt{create-db-cluster} no acepta un array de subredes sino un \textit{subnet group}, que creamos de la forma siguiente:

\begin{lstlisting}[language=bash,basicstyle=\ttfamily\footnotesize]
aws rds create-db-subnet-group \
    --db-subnet-group-name "$DB_SUBNET_GROUP" \
    --db-subnet-group-description "Subnets Aurora EventHub (2 primeras de la VPC)" \
    --subnet-ids $SUBNET_IDS
\end{lstlisting}

\vspace{1em}

\texttt{\$SUBNET\_IDS} contiene los dos identificadores separados por un espacio (por ejemplo, \texttt{subnet-0aaa subnet-0bbb}), y \texttt{--subnet-ids} espera recibir cada uno como un argumento distinto. Por eso la variable va sin comillas, a propósito. Sin comillas, bash la divide por los espacios y la CLI recibe dos argumentos, \texttt{subnet-0aaa} y \texttt{subnet-0bbb}. Con comillas (\texttt{"\$SUBNET\_IDS"}), bash la pasaría como un único argumento, \texttt{"subnet-0aaa subnet-0bbb"}, y la CLI lo rechazaría por no ser un identificador de subred válido.

\subsection{Creación de las instancias}\label{sec:instancias}

El \textit{subnet group} indica en qué subredes, y por tanto en qué zonas de disponibilidad, \emph{pueden} colocarse las instancias del clúster, pero no decide dónde va cada una. Eso se decide al crear cada instancia, y la opción que lo controla no recibe una subred, sino una zona: \texttt{create-db-instance} admite \texttt{--availability-zone}. 

Si no se indica la zona, \textit{Amazon Relational DataBase Service} (RDS, el servicio de AWS que gestiona las bases de datos relacionales) elige una subred cualquiera del conjunto de las subredes del \textit{subnet group}, y nada garantiza que las dos instancias acaben en zonas distintas porque dos subredes podrían estar en la misma zona de disponibilidad. Esto no es el caso de este laboratorio, porque usamos las subredes predefinidas que AWS Academy nos asegura que están definidas en zonas distintas. Sin embargo, esto no tiene por qué ser cierto en un caso general.

Si dos instancias son asignadas a la misma zona, y la zona fallase, esto implicaría que las dos caerían a la vez. Aurora no podría hacer \textit{failover}, que es el mecanismo por el que, si falla la instancia de escritura, se promueve automáticamente una réplica a instancia de escritura para que el servicio continúe (a este respecto, ver la sección~\ref{sec:writer}). Para eso hace falta que la réplica siga en pie. El \textit{failover} solo se produce cuando falla la instancia de escritura, sea cual sea el número de instancias del clúster: es un problema de disponibilidad, porque sin instancia de escritura la aplicación no puede escribir en la base de datos. 

Si falla una réplica de lectura, no hay \textit{failover}, porque ningún papel cambia; parte de las lecturas se quedan sin la réplica que las atendía y recaen en las demás instancias, lo que es una variante del problema de disponibilidad relacionado estrechamente con el de escalabilidad (ver sección~\ref{sec:readers}). Por eso el script indica de forma explícita la zona de cada instancia. Como \texttt{lab-state.json} solo guarda los identificadores de las subredes, primero pregunta a AWS en qué zona está cada una:

\begin{lstlisting}[language=bash,basicstyle=\ttfamily\footnotesize]
# sort_by: el writer va siempre a la primera zona por orden alfabético.
AZS=$(aws ec2 describe-subnets \
    --subnet-ids $SUBNET_IDS \
    --query 'sort_by(Subnets, &AvailabilityZone)[].AvailabilityZone' \
    --output text)
read -r AZ_WRITER AZ_READER <<< "$AZS"
\end{lstlisting}

\vspace{1em}

Aunque no es necesario para separar las instancias, hemos realizado un \texttt{sort\_by} para ordenar el resultado por zona, ya que \texttt{describe-subnets} no garantiza el orden de la respuesta. \texttt{sort\_by} ordena por el nombre de la zona, alfabéticamente (\texttt{us-east-1a} antes que \texttt{us-east-1b}). De esta forma sabemos que la instancia de escritura queda siempre en la primera zona por orden alfabético, y la réplica de lectura, en la segunda.

Con \texttt{--output text}, las dos zonas llegan en una sola línea separadas por un tabulador (por ejemplo, \texttt{us-east-1a \textbackslash t us-east-1b}). \texttt{read} divide esa línea por los espacios y tabuladores: la primera palabra va a \texttt{AZ\_WRITER}, y la segunda, a \texttt{AZ\_READER}. La opción \texttt{-r} hace que \texttt{read} trate las barras invertidas como caracteres normales, y no como escapes. Si solo sale una zona, \texttt{AZ\_READER} queda vacía. En ese caso, o si las dos zonas coinciden, el script termina con error antes de crear ningún recurso.

\subsection{Abrir el puerto 5432 en el grupo de seguridad}\label{sec:puerto}

Para que Spring Boot, en la EC2, pueda conectar con Aurora, el grupo de seguridad tiene que admitir tráfico en el puerto 5432, el de PostgreSQL. Es lo mismo que \texttt{ec2.sh} ya hizo con los puertos 22 (SSH), 80 (HTTP) y 443 (HTTPS), y con el mismo comando, \texttt{\seqsplit{authorize-security-group-ingress}}. Si no se abre, la base de datos no puede recibir conexiones.

La diferencia está en quién puede conectarse. Una regla de entrada indica siempre un puerto y un origen. El origen puede ser un rango de direcciones IP (\texttt{--cidr}) o un grupo de seguridad (\texttt{--source-group}); en el segundo caso, se admite el tráfico que procede de cualquier recurso de ese grupo. Al front y al back queremos que pueda conectarse cualquier cliente, desde cualquier IP; por eso \texttt{ec2.sh} usa \texttt{--cidr 0.0.0.0/0} en las reglas de los puertos 80 y 443, lo que significa que cualquier dirección de Internet puede conectarse a dichos puertos. 

Sin embargo, la situación es diferente en el caso de la base de datos. El único componente que debe conectarse a ella es el back. Por eso \texttt{aurora.sh} usa \texttt{--source-group} con el propio grupo \texttt{eventhub-sg}. Como la instancia EC2 y el clúster Aurora pertenecen al grupo de seguridad, solo la instancia EC2, (donde se ejecuta el back) puede conectar con Aurora; desde Internet, el puerto 5432 sigue cerrado%
\endnote{En un entorno real, Aurora tendría probablemente su propio grupo de seguridad, con una única regla de entrada: el puerto 5432, con origen el grupo de la EC2. No admitiría conexiones en los puertos 22, 80 y 443, aunque Aurora no escuche en ellos. Al compartir el grupo de la EC2, el clúster del laboratorio hereda también esas reglas.}.%

\begin{lstlisting}[language=bash,basicstyle=\ttfamily\footnotesize]
# --source-group: solo los miembros del grupo pueden conectar a Aurora
aws ec2 authorize-security-group-ingress \
    --group-id "$SG_ID" \
    --protocol tcp \
    --port "$DB_PORT" \
    --source-group "$SG_ID"
\end{lstlisting}

\vspace{1em}

\subsection{Crear el clúster con password gestionada por Aurora}\label{sec:crear-cluster}

Al crear el clúster hay que dar una password al usuario \textit{master}, el usuario administrador de la base de datos (el que en PostgreSQL suele llamarse \texttt{postgres}, o \texttt{root} en MySQL). En la versión anterior, con H2, la password estaba en un secreto que creábamos nosotros en Secrets Manager y que Spring Boot leía al arrancar. Aquí no podemos seguir ese camino. Si el secreto lo creásemos nosotros, \texttt{aurora.sh} tendría que leer la password del secreto (\texttt{\seqsplit{get-secret-value}}) y pasársela a RDS en \texttt{\seqsplit{create-db-cluster}} (\texttt{\seqsplit{--master-user-password}}). La password llegaría así al ordenador donde se ejecuta el script, y volvería a ser vulnerable: quedaría en la memoria del proceso, en la línea de comandos y, quizá, en el historial de la \textit{shell}.

Por eso se invierte el orden. El secreto sigue estando en Secrets Manager, pero ya no lo creamos nosotros: lo crea Aurora. Con \texttt{\seqsplit{--manage-master-user-password}}, Aurora genera la password del usuario \textit{master}, la asigna al clúster y la guarda en un secreto de Secrets Manager. Spring Boot lee después ese secreto igual que leía el de H2, pero la password no pasa por el script ni sale de AWS.

\begin{lstlisting}[language=bash,basicstyle=\ttfamily\footnotesize]
# --master-username: en Aurora PostgreSQL no puede ser admin
# --manage-master-user-password: Aurora crea el secreto; la password no pasa por la CLI
# --query: la respuesta ya trae el endpoint y el ARN del secreto
CLUSTER_INFO=$(aws rds create-db-cluster \
    --db-cluster-identifier "$DB_CLUSTER_ID" \
    --engine "$DB_ENGINE" \
    --master-username "$DB_USER" \
    --manage-master-user-password \
    --database-name "$DB_NAME" \
    --port "$DB_PORT" \
    --vpc-security-group-ids "$SG_ID" \
    --db-subnet-group-name "$DB_SUBNET_GROUP" \
    --backup-retention-period 1 \
    --no-deletion-protection \
    --query '[DBCluster.Endpoint, DBCluster.MasterUserSecret.SecretArn]' \
    --output text)
state_set DBClusterIdentifier "$DB_CLUSTER_ID"

# --output text separa los dos valores con un tabulador
read -r DB_HOST SECRET_ARN <<< "$CLUSTER_INFO"
state_set Endpoint "$DB_HOST"
state_set SecretArn "$SECRET_ARN"

aws rds wait db-cluster-available --db-cluster-identifier "$DB_CLUSTER_ID"
\end{lstlisting}

\vspace{1em}

El nombre del secreto que crea Aurora no lo elegimos nosotros. RDS lo forma con el prefijo \texttt{rds!cluster-} seguido de un identificador único que genera AWS, por ejemplo \texttt{\seqsplit{rds!cluster-3f2a9c1e-7b4d-4e8a-9f61-0c5d2b8e4a17}}. El prefijo indica que el secreto lo gestiona RDS y que pertenece a un clúster. Como el nombre no se conoce hasta que se crea el clúster, no puede escribirse de antemano en la configuración de Spring Boot: el script lo obtiene después y lo escribe en \texttt{aurora.properties} (sección~\ref{sec:secreto-endpoint}). El secreto contiene únicamente dos claves, \texttt{username} y \texttt{password}.

\texttt{\seqsplit{--backup-retention-period}} controla las copias de seguridad automáticas. Aurora las hace siempre, sin que haya que programarlas: guarda los cambios del volumen de forma continua, y con ellos puede restaurar el clúster tal como estaba en cualquier instante del periodo de retención. La opción fija la duración de ese periodo, en días, entre 1 y 35; las copias automáticas no pueden desactivarse. El script usa 1, el mínimo, porque en el laboratorio no hace falta recuperar datos antiguos y así se ocupa menos almacenamiento de copias. En producción, el valor depende de cuánto tiempo atrás se quiera poder volver, por ejemplo tras un borrado accidental.

\texttt{create-db-cluster} devuelve el objeto \texttt{DBCluster} completo aunque el clúster siga en estado \texttt{creating}. Esa respuesta ya incluye el \textit{endpoint} de escritura y el ARN del secreto, de modo que no hace falta un \texttt{describe-db-clusters} posterior. El nombre DNS del \textit{endpoint} se fija al crear el clúster y no cambia durante su vida: en un failover cambia la instancia a la que apunta, no el nombre. Como el secreto gestionado no incluye los datos de conexión, el \textit{endpoint} se guarda aquí para escribirlo después en \texttt{aurora.properties}.

El script no pasa \texttt{--storage-encrypted}, que activa el cifrado en reposo. En producción convendría añadirlo, porque el cifrado solo puede activarse al crear el clúster.

\texttt{aws rds wait db-cluster-available} bloquea hasta que el estado del clúster es \texttt{available}. Sin esa espera, \texttt{create-db-instance} fallaría porque el clúster aún no está listo. Puede tardar varios minutos.

\subsection{Instancias en dos zonas}\label{sec:dos-zonas}

Un clúster Aurora sin instancias no acepta conexiones JDBC, y para que funcione bastaría con una. El script crea dos, cada una en una de las zonas obtenidas en la sección~\ref{sec:instancias}, por dos motivos:

\begin{itemize}
    \item \textbf{Disponibilidad}: si falla la instancia de escritura, o toda su zona, Aurora promueve la réplica a instancia de escritura (ver sección~\ref{sec:writer}).
    \item \textbf{Escalabilidad}: la réplica atiende las lecturas de la aplicación, descargando de trabajo a la instancia de escritura (secciones~\ref{sec:endpoints} y~\ref{sec:escalabilidad}). Si la aplicación tuviera que atender a más usuarios, podrían añadirse más réplicas de lectura, hasta quince por clúster, con el mismo \texttt{\seqsplit{create-db-instance}}. Las réplicas solo reparten las lecturas: todas las escrituras son procesadas por la única instancia de escritura.
\end{itemize}

Primero se crea la de escritura:

\begin{lstlisting}[language=bash,basicstyle=\ttfamily\footnotesize]
# --availability-zone: zona de la primera subnet
# --db-instance-class: clase mínima habitual de Aurora PostgreSQL
# --no-publicly-accessible: solo se alcanza desde la VPC
aws rds create-db-instance \
    --db-instance-identifier "$DB_WRITER_ID" \
    --db-cluster-identifier "$DB_CLUSTER_ID" \
    --engine "$DB_ENGINE" \
    --db-instance-class db.t3.medium \
    --availability-zone "$AZ_WRITER" \
    --no-publicly-accessible
state_set_array DBInstanceIdentifiers "$DB_WRITER_ID"

aws rds wait db-instance-available --db-instance-identifier "$DB_WRITER_ID"
\end{lstlisting}

\vspace{1em}

Después, la réplica de lectura:

\begin{lstlisting}[language=bash,basicstyle=\ttfamily\footnotesize]
# --availability-zone: zona de la segunda subnet, distinta de la del writer
# --promotion-tier 0: prioridad máxima para sustituir al writer si falla
aws rds create-db-instance \
    --db-instance-identifier "$DB_READER_ID" \
    --db-cluster-identifier "$DB_CLUSTER_ID" \
    --engine "$DB_ENGINE" \
    --db-instance-class db.t3.medium \
    --availability-zone "$AZ_READER" \
    --promotion-tier 0 \
    --no-publicly-accessible
state_set_array DBInstanceIdentifiers "$DB_WRITER_ID" "$DB_READER_ID"

aws rds wait db-instance-available --db-instance-identifier "$DB_READER_ID"
\end{lstlisting}

\vspace{1em}

La primera instancia que se crea en un clúster Aurora es la de escritura; las siguientes son réplicas. Por eso el script espera a que la de escritura esté \texttt{available} antes de crear la réplica: si las crease a la vez, cualquiera de las dos podría acabar como instancia de escritura, y no sabríamos en qué zona queda. El precio es un despliegue más largo, porque las dos instancias se crean secuencialmente.

\texttt{--promotion-tier} asigna a la instancia un nivel de prioridad para el \textit{failover}: un número entre 0 y 15, que por defecto es 1. Cuanto más bajo es el número, más prioridad tiene. Si falla la instancia de escritura, Aurora promueve la réplica con el número más bajo. Varias réplicas pueden tener el mismo nivel; en caso de empate, Aurora elige la de mayor tamaño (clase de instancia) y, si también coinciden, una cualquiera de ellas. El 15 no tiene relación con el número máximo de réplicas: es solo el nivel más bajo de prioridad. Con una sola réplica, como en \proyecto, no hay elección posible, pero el script indica el 0 para que quede explícito que esa réplica es la candidata a sustituir a la instancia de escritura.

\texttt{--db-instance-class} fija el tamaño de la instancia, es decir, su CPU y su memoria. El script usa \texttt{\seqsplit{db.t3.medium}} (2 vCPU y 4 GiB de RAM) porque es la clase más pequeña de las disponibles en AWS Academy para Aurora PostgreSQL. Es una clase de rendimiento variable (\textit{burstable}), pensada para desarrollo y pruebas con poca carga, suficiente para \proyecto. Elegir la más pequeña importa porque una instancia de Aurora es un recurso relativamente caro: está en marcha de forma permanente, reserva su CPU y su memoria para el motor de base de datos, y se factura por horas aunque no reciba ninguna consulta. El clúster de \proyecto~tiene dos, así que el coste se duplica. En AWS Academy ese coste se descuenta del crédito del laboratorio mientras el clúster exista, por eso conviene borrarlo al terminar con \texttt{aurora.sh delete}.

\texttt{\seqsplit{--no-publicly-accessible}} hace que la instancia tenga solo una dirección IP privada, dentro de la VPC, y ninguna IP pública. Es una protección distinta de la del grupo de seguridad: el grupo decide quién puede conectarse al puerto 5432, pero la IP privada hace que, desde fuera de la VPC, no haya siquiera una dirección a la que intentarlo. Spring Boot no lo nota, porque la EC2 está en la misma VPC y llega al clúster por su IP privada. La opción se indica de forma explícita porque, si se omite, RDS hace pública la instancia cuando sus subredes tienen salida a Internet, como ocurre con las de la VPC por defecto.

\subsection{Secreto, \textit{endpoint} y \texttt{aurora.properties}}\label{sec:secreto-endpoint}

Cuando las dos instancias están \texttt{available}, el clúster ya funciona, pero Spring Boot todavía no sabe cómo llegar a él. Necesita tres cosas que solo se conocen después de crear el clúster, y que cambian en cada despliegue:

\begin{itemize}
    \item \textbf{El nombre del secreto}, para leer de Secrets Manager el usuario y la password del usuario \textit{master}.
    \item \textbf{El \textit{endpoint} del clúster}, es decir, el nombre DNS al que se conecta el driver JDBC.
    \item \textbf{El puerto y el nombre de la base de datos}, que completan la URL JDBC.
\end{itemize}

El script ya tiene todo. El \textit{endpoint} y el ARN del secreto se obtuvieron de la respuesta de \texttt{create-db-cluster}. El puerto y el nombre de la base son constantes del script. Spring Boot podría leer el secreto directamente con su ARN, pero, por sencillez, el script le pasa solo el nombre del secreto (\texttt{rds!cluster-}\ldots), que es más corto y legible. Para obtenerlo a partir del ARN hace una consulta más a Secrets Manager.

\paragraph{Obtener el nombre del secreto.}

\begin{lstlisting}[language=bash,basicstyle=\ttfamily\footnotesize]
# El nombre del secreto lo asigna RDS (rds!cluster-...). describe-secret no devuelve la password.
SECRET_NAME=$(aws secretsmanager describe-secret \
    --secret-id "$SECRET_ARN" \
    --query Name \
    --output text)
\end{lstlisting}

\vspace{1em}

Secrets Manager distingue entre los \textbf{metadatos} de un secreto (su nombre, su ARN, su descripción, sus fechas o su configuración de rotación) y su \textbf{contenido} (el \texttt{SecretString}, que en este caso es un JSON con \texttt{username} y \texttt{password}). \texttt{describe-secret} devuelve solo los metadatos; el contenido únicamente lo devuelve \texttt{\seqsplit{get-secret-value}}, que el script no llama nunca. Así, el script averigua cómo se llama el secreto sin llegar a ver la password.

\texttt{--secret-id} identifica el secreto; se le pasa el ARN, que es lo que el script tiene. \texttt{--query Name} se queda solo con el campo \texttt{Name} de la respuesta, y \texttt{--output text} lo imprime sin comillas, listo para guardarlo en \texttt{SECRET\_NAME}. Como el secreto lo gestiona RDS, su vida está ligada a la del clúster: al borrar el clúster con \texttt{aurora.sh delete}, RDS elimina también el secreto, así que no quedan secretos huérfanos en Secrets Manager.

Finalmente, los cuatro datos se escriben en el fichero \texttt{\seqsplit{aurora.properties}} para que puedan ser usados por el back. El resultado tiene este aspecto:

\begin{lstlisting}[language=properties,basicstyle=\ttfamily\footnotesize]
# Generado por aws-scripts/aurora.sh. No editar a mano.
# El secreto lo crea Aurora y aporta username y password. Spring lo importa al arrancar.
# aurora.host, aurora.port y aurora.dbname no son secretos: construyen la URL JDBC.
spring.config.import=aws-secretsmanager:rds!cluster-xxxxxxxxxxxx
aurora.host=eventhub-cluster-aurora.cluster-xxxxxxxxxxxx.us-east-1.rds.amazonaws.com
aurora.port=5432
aurora.dbname=eventhub
\end{lstlisting}

\section{\textit{Endpoints} y reparto de lecturas y escrituras}\label{sec:endpoints}

Las aplicaciones no se conectan a una instancia concreta, sino a un nombre DNS (un \textit{endpoint}). Un clúster Aurora ofrece tres tipos de endpoints:

\begin{itemize}
    \item \textbf{\textit{Endpoint} del clúster}: \texttt{\seqsplit{eventhub-cluster-aurora.cluster-}}\ldots\texttt{\seqsplit{.rds.amazonaws.com}}. Apunta siempre a la instancia de escritura actual; tras un \textit{failover}, pasa a apuntar a la instancia promovida.
    \item \textbf{\textit{Endpoint} de lectura}: \texttt{\seqsplit{eventhub-cluster-aurora.cluster-ro-}}\ldots\texttt{\seqsplit{.rds.amazonaws.com}}. Reparte las conexiones entre las réplicas de lectura.
    \item \textbf{\textit{Endpoints} de instancia}: \texttt{\seqsplit{eventhub-aurora-writer.}}\ldots{} y \texttt{\seqsplit{eventhub-aurora-reader.}}\ldots{} Apuntan a una instancia concreta, sea cual sea su papel en cada momento.
\end{itemize}

\paragraph{Usar los \textit{endpoints} directamente.} Cada \textit{endpoint} puede usarse por separado, con el driver estándar de PostgreSQL y una URL \texttt{\seqsplit{jdbc:postgresql://}}\ldots{}, pero entonces es la aplicación la que tiene que decidir a cuál se conecta, esto es, por qué endpoint envía los comandos SQL de lectura y por cual los de escritura. Por el contrario, si la aplicación usa solo endpoint del clúster, todos los comandos SQL, de lectura y de escritura, van a la instancia de escritura, y las réplicas no se aprovechan. 

Además, el \textit{endpoint} de lectura no es un balanceador de carga. Es un nombre DNS que, cada vez que se resuelve, devuelve la dirección de una de las réplicas, rotando entre ellas (\textit{round-robin}) con un \textit{Time To Live} (TTL) de 5 segundos. Se reparten \textbf{conexiones}, no consultas: una vez abierta, todas las consultas de una conexión van a la misma réplica. El reparto tampoco tiene en cuenta la carga de cada réplica. 

\paragraph{El driver JDBC de AWS.} \proyecto~usa otra opción: el driver JDBC de AWS, que reparte las lecturas y las escrituras entre las instancias sin que la aplicación tenga que elegir a qué \textit{endpoint} se conecta. Le basta una única URL, la del \textit{endpoint} del clúster. 

El driver no analiza las consultas: lo que ve es que Spring marca la conexión como de solo lectura al empezar una transacción \texttt{@Transactional(readOnly = true)}, y la desmarca al terminar. Con esa señal, el driver envía la transacción a una réplica; el resto de transacciones van a la instancia de escritura. Por ejemplo, en \proyecto~. el\texttt{EventoService} marca la búsqueda de eventos:

\begin{lstlisting}[language=Java,basicstyle=\ttfamily\footnotesize]
@Transactional(readOnly = true)
public List<EventoResponseDTO> buscarEventos(String nombre) {
    ...
}
\end{lstlisting}

\vspace{1em}

\noindent de modo que esa consulta va a \texttt{eventhub-aurora-reader} sin que el código tenga que saber nada de Aurora.

Como el driver se conecta a los \textit{endpoints} de instancia usando la topología, no depende del \textit{endpoint} de lectura ni de su reparto por DNS.

\section{Disponibilidad}\label{sec:disponibilidad}

Esta sección describe qué papel desempeña cada tipo de instancia del clúster, qué ocurre cuando falla y cómo afecta al tiempo de respuesta.

\subsection{La Instancia de Escritura}\label{sec:writer}

La instancia de escritura acepta escrituras (\texttt{INSERT}, \texttt{UPDATE}, \texttt{DELETE} y cambios de esquema) y también lecturas. Cada cambio que confirma se escribe en el volumen de almacenamiento compartido, con copias en tres zonas. Un clúster Aurora PostgreSQL tiene \textbf{una única instancia de escritura}: las réplicas añaden capacidad de lectura, no de escritura.

\subsubsection{Si falla la instancia de escritura}

\begin{itemize}
    \item \textbf{\textit{Failover}}. Aurora promueve una réplica a instancia de escritura (sección~\ref{sec:instancias}); si hay varias, elige según su \texttt{--promotion-tier} (sección~\ref{sec:dos-zonas}). El proceso suele terminar en menos de 60 segundos, y a menudo en menos de 30. Con una sola instancia no habría réplica que promover: Aurora tendría que crear otra, lo que suele llevar menos de 10 minutos, pero sin garantía.
    \item \textbf{Los papeles cambian, los nombres no}. Tras un \textit{failover}, \texttt{\seqsplit{eventhub-aurora-reader}} pasa a ser la instancia de escritura, y \texttt{\seqsplit{eventhub-aurora-writer}}, cuando se recupera, vuelve como réplica. El \textit{endpoint} del clúster sigue teniendo el mismo nombre, pero apunta a la nueva instancia de escritura.
    \item \textbf{Si cae toda la zona}, la réplica está en la otra y el \textit{failover} funciona igual. Por eso el script coloca las dos instancias en zonas distintas (sección~\ref{sec:instancias}).
    \item \textbf{Los datos no se pierden}: están en la capa de almacenamiento, replicados en tres zonas, independientemente de dónde estén las instancias.
    \item \textbf{El \textit{failover} no es transparente para la aplicación}. Durante esos segundos, las conexiones abiertas fallan y las transacciones en curso se abortan. El \textit{endpoint} del clúster es un nombre DNS que pasa a apuntar a la nueva instancia de escritura, con un TTL de 5 segundos. El driver JDBC de AWS no espera a que caduque: el \textit{plugin} \texttt{failover2} detecta la nueva instancia de escritura consultando la topología del clúster, y \texttt{efm2} (\textit{enhanced failure monitoring}) vigila la instancia a la que está conectada cada conexión para detectar antes que ha caído. El \textit{pool} de conexiones de Spring Boot (HikariCP) descarta las conexiones rotas y abre otras nuevas.
    \item \textbf{La caché de la instancia promovida}. La réplica tiene en memoria las páginas que ha usado para sus lecturas, no las de la instancia de escritura, así que tras el \textit{failover} las primeras consultas pueden ser más lentas. Aurora PostgreSQL permite que la réplica mantenga su caché al día con la de la instancia de escritura (\textit{cluster cache management}). Se activa con un parámetro del clúster.
\end{itemize}

Todo lo anterior protege frente a la caída de una zona, no de una región entera. Para eso, Aurora ofrece \textit{Global Database}, que replica el clúster en otras regiones de AWS con un retraso habitualmente inferior a un segundo y permite promover una de ellas si falla la región principal. \proyecto~no la usa.

Aurora tampoco protege al resto de la aplicación. Hay una sola instancia EC2, en una zona que elige AWS. Si cae esa zona, la aplicación deja de responder aunque Aurora siga funcionando. La alta disponibilidad completa exigiría varias EC2 en distintas zonas detrás de un balanceador de carga.

\subsubsection{Tiempo de respuesta}

\begin{itemize}
    \item \textbf{Las escrituras no se ralentizan por tener una réplica}. Cada escritura ya se guarda en la capa de almacenamiento, con copias en tres zonas. La réplica no participa en ese camino: solo recibe de la instancia de escritura, de forma asíncrona, los cambios que necesita para actualizar su caché.
    \item \textbf{La latencia depende de la zona de la EC2}. Una consulta entre la EC2 y la instancia de escritura en la misma zona tarda del orden de décimas de milisegundo en ir y volver; entre zonas distintas, del orden de un milisegundo. Hibernate lanza varias consultas por petición, así que la diferencia es de milisegundos. \texttt{ec2.sh} no elige la zona de la EC2, y tras un \textit{failover} la instancia de escritura cambia de zona.
\end{itemize}

\subsection{Las Réplicas de Lectura}\label{sec:readers}

Una réplica de lectura solo atiende lecturas: un intento de escritura falla con un error del tipo \texttt{cannot execute INSERT in a read-only transaction}. No escribe en el almacenamiento; la instancia de escritura le envía el flujo de cambios para que actualice las páginas que tiene en memoria.

\begin{itemize}
    \item \textbf{Las réplicas van por detrás}. Por ese envío de cambios, una réplica va normalmente menos de 100 milisegundos por detrás de la instancia de escritura. Si un usuario compra una entrada e inmediatamente consulta sus compras en una réplica, puede no ver la compra. Las lecturas que deben ver lo recién escrito tienen que ir a la instancia de escritura, es decir, no marcarse como de solo lectura (sección~\ref{sec:endpoints}).
    \item \textbf{Si falla una réplica}, no hay \textit{failover}: la instancia de escritura sigue atendiendo escrituras y los datos están en la capa de almacenamiento, así que la disponibilidad no se ve afectada. Aurora reinicia la réplica y, mientras tanto, las conexiones que el driver tenía abiertas con ella fallan; en la siguiente transacción de solo lectura, el driver abre una conexión con otra réplica, o con la instancia de escritura si no queda ninguna. Lo que se pierde es escalabilidad: hasta que la réplica se recupera, sus lecturas recaen en las demás instancias. En \proyecto, con una sola réplica, todas las lecturas vuelven a la instancia de escritura.
\end{itemize}

\section{Escalabilidad}\label{sec:escalabilidad}

En Aurora, las lecturas y las escrituras escalan de forma distinta. Las lecturas se reparten entre réplicas, que pueden añadirse y quitarse según la carga. Las escrituras las atiende siempre una única instancia, así que solo escalan haciéndola más potente o quitándole trabajo.

\subsection{Escalar las lecturas}

La mayoría de las aplicaciones leen mucho más de lo que escriben; en \proyecto, por ejemplo, listar eventos es mucho más frecuente que comprar. Por eso la forma habitual de escalar un clúster Aurora es añadir réplicas de lectura. Como una réplica no es una copia de los datos, sino una instancia más conectada al volumen compartido (sección~\ref{sec:aurora}), añadirla o quitarla es una operación ligera. Puede hacerse a mano o de forma automática, y en los dos casos la réplica se integra en el clúster del mismo modo.

\paragraph{Añadir y quitar réplicas a mano.} Las réplicas pueden crearse y borrarse libremente mientras el clúster está en marcha. Una réplica nueva se crea con \texttt{\seqsplit{create-db-instance}} indicando el clúster al que pertenece (sección~\ref{sec:dos-zonas}), y se borra con \texttt{\seqsplit{delete-db-instance}}. Al crearla conviene elegir su zona con \texttt{\seqsplit{--availability-zone}}, por la misma razón que en la sección~\ref{sec:instancias}: que las instancias no se concentren en una sola zona. Crearla no obliga a copiar datos: tarda los minutos que cuesta arrancar la instancia, sea cual sea el tamaño de la base de datos. Ni crear ni borrar una réplica afecta a la instancia de escritura. Borrar la instancia de escritura es distinto: provoca un \textit{failover} (sección~\ref{sec:writer}).

\paragraph{Autoescalado.} \textit{Aurora Auto Scaling} hace lo mismo de forma automática. Se le indica un número mínimo y máximo de réplicas y una métrica que vigilar, por ejemplo la CPU media de las réplicas o el número de conexiones, y añade réplicas cuando la carga sube y las quita cuando baja. Así el número de réplicas, y con él el coste, sigue a la demanda sin intervención manual.

\paragraph{Cómo se integran las réplicas.} Tanto si se añaden a mano como con autoescalado, la integración es casi transparente:

\begin{itemize}
    \item \textbf{En el clúster} no hay que configurar nada. El \textit{endpoint} de lectura (sección~\ref{sec:endpoints}) incluye la réplica nueva en cuanto está \texttt{available}, y deja de devolver la que se borra.
    \item \textbf{En la aplicación}, con el driver JDBC de AWS, tampoco hay que cambiar la configuración ni reiniciar nada. El driver descubre las réplicas nuevas cuando refresca la topología del clúster, lo que hace periódicamente (por defecto, cada 30 segundos). Solo las usan las conexiones con réplica que se abren después: las conexiones del \textit{pool} que ya tenían una réplica asignada siguen con ella hasta que se renuevan (sección~\ref{sec:endpoints}).
    \item \textbf{Al borrar una réplica}, la transacción de solo lectura que estuviera en curso en ella falla, y la aplicación recibe el error. En la siguiente, el driver abre una conexión con otra réplica, o con la instancia de escritura si no queda ninguna. Es lo mismo que ocurre cuando una réplica falla (sección~\ref{sec:readers}).
\end{itemize}

\paragraph{Más allá de quince réplicas.} Un clúster admite como máximo quince réplicas (sección~\ref{sec:aurora}), dieciséis instancias contando la de escritura. Es un límite fijo de Aurora, y también el máximo hasta el que puede crecer \textit{Aurora Auto Scaling}. Cuando no basta, hay tres caminos:

\begin{itemize}
    \item \textbf{Réplicas más grandes.} Cambiar la clase de las réplicas por una mayor (\texttt{\seqsplit{modify-db-instance}}) hace que cada una atienda más lecturas.
    \item \textbf{Más clústeres que reciben los cambios.} \textit{Aurora Global Database} añade clústeres secundarios, de solo lectura, en otras regiones de AWS. Cada uno tiene sus propias réplicas y recibe los cambios del clúster principal con un retraso habitualmente inferior a un segundo. La aplicación tiene que conectarse a ellos: su driver no los ve como parte del clúster principal.
    \item \textbf{Quitar lecturas a la base de datos.} Una caché delante de Aurora atiende las lecturas repetidas sin llegar a la base de datos. Es un elemento nuevo en la arquitectura, que se trata en otro documento.
\end{itemize}

En la práctica, el límite suele aparecer antes por el coste: cada réplica es una instancia que se paga por horas aunque no reciba consultas.

\subsection{Escalar las escrituras}

Todas las escrituras van a la única instancia de escritura, así que las réplicas no ayudan con ellas. Antes de actuar conviene medir dónde está el cuello de botella; a menudo no es la base de datos, sino la aplicación. Si lo son las escrituras, hay tres opciones que no añaden ningún elemento a la arquitectura:

\begin{itemize}
    \item \textbf{Escalado vertical}: cambiar la clase de la instancia (\texttt{\seqsplit{aws\ rds\ modify-db-instance\ --db-instance-class}}) por una mayor. Las clases de memoria optimizada (\texttt{db.r7g.*}, por ejemplo) llegan a decenas de vCPU y cientos de GiB de RAM. El cambio reinicia la instancia; para evitar la interrupción, se cambia primero la réplica y después se provoca un \textit{failover} hacia ella.
    \item \textbf{Aurora Serverless v2}: las instancias \texttt{db.serverless} ajustan su capacidad a la carga, entre un mínimo y un máximo de ACU (\textit{Aurora Capacity Units}; una ACU equivale a unos 2 GiB de RAM con su CPU y red correspondientes), sin reiniciarse. Encaja con picos de demanda, como la apertura de ventas de un concierto.
    \item \textbf{Optimizar las consultas}: añadir índices, evitar el problema de las N+1 consultas de Hibernate o agrupar escrituras en lotes. A menudo es lo que más rendimiento aporta con menos coste.
\end{itemize}

Cuando ni así basta, hay que introducir elementos nuevos en la arquitectura (un \textit{proxy} de conexiones, colas o el reparto de los datos entre varias bases de datos), que quedan fuera de este documento.

\section{Conclusión}

Amazon Aurora ofrece una arquitectura relacional de clase empresarial: almacenamiento distribuido y replicado, cómputo independiente del disco y failover automático cuando el clúster dispone de réplicas. En \proyecto~esa arquitectura se concreta en un clúster Aurora PostgreSQL con una instancia de escritura y una réplica en zonas de disponibilidad distintas, cuya contraseña genera el propio RDS al crear el clúster (\texttt{--manage-master-user-password}), de modo que la clave no transita por la CLI. El nombre del secreto lo asigna AWS; \texttt{aurora.sh} lo deja, junto con el \textit{endpoint} del clúster, en \texttt{aurora.properties} para que Spring Boot monte el \textit{datasource} sin credenciales en el repositorio. El grupo de seguridad de la EC2, transmitido por \texttt{lab-state.json}, aísla el puerto 5432 al interior de la VPC. El resultado es un entorno alineado con las prácticas de secretos, red y persistencia exigibles en una aplicación desplegada en AWS, en el que la base de datos resiste la caída de una instancia o de una zona completa. El siguiente paso hacia la alta disponibilidad no está en Aurora, sino en la EC2, que sigue siendo un punto único de fallo.

\bibliographystyle{plainnat}
\bibliography{main}

\theendnotes

\end{document}
